\documentclass[10pt,a4paper]{article}
\usepackage[T1]{fontenc}
\usepackage[margin=2cm]{geometry}
\usepackage{isabelle,isabellesym}
\usepackage{amsmath,amssymb}
\usepackage{pdfsetup}

\urlstyle{rm}
\isabellestyle{it}

\begin{document}

\title{Migotti's Theorem on the Coefficients of the Cyclotomic Polynomial \(\Phi_{pq}\)}
\author{Carles Mar\'in}
\maketitle

\begin{abstract}
Migotti (1883)~\cite{Migotti1883} proved that for distinct primes $p$ and $q$ every coefficient of the
cyclotomic polynomial $\Phi_{pq}$ is $-1$, $0$ or $1$. We formalize this for the polynomial
$\Phi_{pq} = (1-X^{pq})(1-X)/((1-X^p)(1-X^q))$, which equals the cyclotomic polynomial when $p\neq q$ are primes
(the identification is classical and is not formalized here). More generally, for coprime $p,q>0$ the quotient is
exact and has coefficients in $\{-1,0,1\}$. The proof writes $\Phi_{pq}=1-(1-X)B$, where $B$ is a sum of distinct
monomials $X^{pa+qb-pq}$, using that $(a,b)\mapsto (pa+qb)\bmod pq$ is a bijection.
\end{abstract}

\paragraph{Acknowledgement.} The formalization and the text of this document were prepared with the
assistance of generative AI systems, and checked by the author.

\tableofcontents

\parindent 0pt\parskip 0.5ex

\input{session}

\bibliographystyle{abbrv}
\bibliography{root}

\end{document}
